\begin{tomtat}
	\begin{minipage}[t]{0.3\textwidth}
		\vspace{0pt}
		\begin{tomtat}
			\begin{center}
				\textbf{ĐỀ 09}
			\end{center}
			
		\end{tomtat}
	\end{minipage}
	\hfill
	\begin{minipage}[t]{0.5\textwidth}
		\vspace{0pt}
		\begin{tomtat}
			\begin{center}
				\includegraphics[scale=0.08]{images/9}
			\end{center}
		\end{tomtat}
	\end{minipage}
\end{tomtat}
\textbf{I. PHẦN TRẮC NGHIỆM}
\begin{ex}
	Kết quả của phép khai căn $\sqrt{(\sqrt{11}-4)^2}$ là:
	\choice{$4-\sqrt{11}$}{$-4-\sqrt{11}$}{$\sqrt{11}-4$}{$\sqrt{11}+4$}
	\loigiai{
		Ta có $\sqrt{A^2} = |A|$.
		Do đó, $\sqrt{(\sqrt{11}-4)^2} = |\sqrt{11}-4|$.
		Để xác định giá trị tuyệt đối, ta so sánh $\sqrt{11}$ và $4$.
		Ta có $11 < 16 \Rightarrow \sqrt{11} < \sqrt{16} \Rightarrow \sqrt{11} < 4$.
		Vì $\sqrt{11} < 4$ nên $\sqrt{11}-4 < 0$.
		Do đó, $|\sqrt{11}-4| = -(\sqrt{11}-4) = 4-\sqrt{11}$.
		\\ \textbf{ĐÁP ÁN A}
	}
\end{ex}

\begin{ex}
	Nghiệm của phương trình $2x-3=15$ là:
	\choice{9}{-9}{6}{-6}
	\loigiai{
		Ta giải phương trình:
		\begin{align*}
			2x-3 &= 15 \\
			2x &= 15+3 \\
			2x &= 18 \\
			x &= \frac{18}{2} \\
			x &= 9
		\end{align*}
		Vậy nghiệm của phương trình là $x=9$.
		\\ \textbf{ĐÁP ÁN A}
	}
\end{ex}

\begin{ex}
	Cho tam giác ABC vuông tại A, có AB=3 cm, BC=5 cm thì tanC bằng:
	\choice{$\frac{3}{5}$}{$\frac{3}{4}$}{$\frac{4}{3}$}{$\frac{5}{3}$}
	\loigiai{
		Trong tam giác ABC vuông tại A, theo định lý Pitago, ta có:
		$AC^2 = BC^2 - AB^2$
		$AC^2 = 5^2 - 3^2$
		$AC^2 = 25 - 9$
		$AC^2 = 16$
		$AC = \sqrt{16} = 4 \text{ cm}$.
		
		Theo định nghĩa tỉ số lượng giác tang trong tam giác vuông, ta có:
		$$ \tan C = \frac{\text{cạnh đối}}{\text{cạnh kề}} = \frac{AB}{AC} $$
		Với $AB=3 \text{ cm}$ và $AC=4 \text{ cm}$, ta có:
		$$ \tan C = \frac{3}{4} $$
		\\ \textbf{ĐÁP ÁN B}
	}
\end{ex}

\begin{ex}
	Nghiệm của bất phương trình $-3x-6 \le 0$ là:
	\choice{$x \le 2$}{$x \ge 2$}{$x \ge -2$}{$x \le -2$}
	\loigiai{
		Ta giải bất phương trình:
		\begin{align*}
			-3x-6 &\le 0 \\
			-3x &\le 6
		\end{align*}
		Khi chia cả hai vế của bất phương trình cho một số âm (ở đây là -3), ta phải đổi chiều bất đẳng thức:
		\begin{align*}
			x &\ge \frac{6}{-3} \\
			x &\ge -2
		\end{align*}
		Vậy nghiệm của bất phương trình là $x \ge -2$.
		\\ \textbf{ĐÁP ÁN C}
	}
\end{ex}

\begin{ex}
	Đồ thị hàm số $y=ax^2 (a \ne 0)$ đi qua điểm A(-2;4) thì a bằng:
	\choice{$-\frac{1}{2}$}{$\frac{1}{2}$}{-1}{1}
	\loigiai{
		Đồ thị hàm số $y=ax^2$ đi qua điểm A(-2;4) có nghĩa là khi $x=-2$ thì $y=4$.
		Thay $x=-2$ và $y=4$ vào phương trình hàm số, ta được:
		\begin{align*}
			4 &= a(-2)^2 \\
			4 &= a \cdot 4 \\
			a &= \frac{4}{4} \\
			a &= 1
		\end{align*}
		Vậy giá trị của $a$ là 1.
		\\ \textbf{ĐÁP ÁN D}
	}
\end{ex}

\begin{ex}
	Gieo con xúc xắc 50 lần cho kết quả như sau:
	\begin{center}
		\begin{tabular}{|l|c|c|c|c|c|c|}
			\hline
			Số chấm xuất hiện & 1 & 2 & 3 & 4 & 5 & 6 \\
			\hline
			Tần số & 7 & 5 & ? & 6 & 8 & 13 \\
			\hline
		\end{tabular}
	\end{center}
	Tần số xuất hiện của mặt 3 chấm là:
	\choice{8}{10}{11}{13}
	\loigiai{
		Tổng số lần gieo là 50.
		Tổng tần số của các mặt đã biết là: $7 (1 \text{ chấm}) + 5 (2 \text{ chấm}) + 6 (4 \text{ chấm}) + 8 (5 \text{ chấm}) + 13 (6 \text{ chấm}) = 39$.
		Tần số xuất hiện của mặt 3 chấm là:
		$$ 50 - 39 = 11 $$
		Vậy tần số xuất hiện của mặt 3 chấm là 11.
		\\ \textbf{ĐÁP ÁN C}
	}
\end{ex}
\begin{ex}
	Một hộp có 30 thẻ cùng loại, mỗi thẻ được ghi một trong các số 1; 2; 3; 4; 5; ...; 29; 30; hai thẻ khác nhau thì ghi số khác nhau. Rút ngẫu nhiên một thẻ trong hộp. Xác suất của biến cố “Số xuất hiện trên thẻ được rút ra là số chia hết cho cả 2 và 5” là:
	\choice{$\frac{2}{3}$}{$\frac{1}{10}$}{$\frac{4}{5}$}{$\frac{5}{6}$}
	\loigiai{
		Tập hợp tất cả các kết quả có thể xảy ra khi rút ngẫu nhiên một thẻ từ hộp là các số từ 1 đến 30.
		Không gian mẫu $\Omega = \{1, 2, 3, ..., 30\}$.
		Số phần tử của không gian mẫu là $n(\Omega) = 30$.
		
		Gọi A là biến cố “Số xuất hiện trên thẻ được rút ra là số chia hết cho cả 2 và 5”.
		Một số chia hết cho cả 2 và 5 thì chia hết cho Bội chung nhỏ nhất của (2,5), tức là chia hết cho 10.
		Tập hợp các kết quả thuận lợi cho biến cố A (các số từ 1 đến 30 chia hết cho 10) là:
		$A = \{10, 20, 30\}$.
		Số phần tử của biến cố A là $n(A) = 3$.
		
		Xác suất của biến cố A là:
		$$ P(A) = \frac{n(A)}{n(\Omega)} = \frac{3}{30} = \frac{1}{10} $$
		\\ \textbf{ĐÁP ÁN B}
	}
\end{ex}

\begin{ex}
	Một hình nón có độ dài đường sinh bằng 6 cm, bán kính đáy bằng 4 cm. Diện tích xung quanh của hình nón bằng:
	\choice{$48\pi \text{ (cm}^2\text{)}$}{$24\pi \text{ (cm}^2\text{)}$}{$12\pi \text{ (cm}^2\text{)}$}{$96\pi \text{ (cm}^2\text{)}$}
	\loigiai{
		Diện tích xung quanh của hình nón được tính bằng công thức: $S_{xq} = \pi Rl$.
		Trong đó:
		$R$ là bán kính đáy của hình nón.
		$l$ là độ dài đường sinh của hình nón.
		Theo đề bài, $R = 4 \text{ cm}$ và $l = 6 \text{ cm}$.
		\begin{align*}
			S_{xq} &= \pi \cdot R \cdot l \\
			&= \pi \cdot 4 \cdot 6 \\
			&= 24\pi \text{ (cm}^2\text{)}
		\end{align*}
		\\ \textbf{ĐÁP ÁN B}
	}
\end{ex}
\newpage
\textbf{II. PHẦN TỰ LUẬN}
\begin{ex} % Câu 9
	Rút gọn biểu thức sau: $P = \left(\frac{2}{\sqrt{x}+1} - \frac{\sqrt{x}-2}{x+\sqrt{x}}\right) : \frac{2}{\sqrt{x}+1}$ với $x>0$.
	\loigiai{
		Với $x>0$, ta có $x+\sqrt{x} = \sqrt{x}(\sqrt{x}+1)$.
		Điều kiện $x>0$ đảm bảo $\sqrt{x}$ xác định và $\sqrt{x}+1 \ne 0$, $x+\sqrt{x} \ne 0$.
		\begin{align*}
			P &= \left(\frac{2}{\sqrt{x}+1} - \frac{\sqrt{x}-2}{\sqrt{x}(\sqrt{x}+1)}\right) : \frac{2}{\sqrt{x}+1} \\
			&= \left(\frac{2\sqrt{x}}{{\sqrt{x}(\sqrt{x}+1)}} - \frac{\sqrt{x}-2}{\sqrt{x}(\sqrt{x}+1)}\right) \cdot \frac{\sqrt{x}+1}{2} \\
			&= \frac{2\sqrt{x} - (\sqrt{x}-2)}{\sqrt{x}(\sqrt{x}+1)} \cdot \frac{\sqrt{x}+1}{2} \\
			&= \frac{2\sqrt{x} - \sqrt{x} + 2}{\sqrt{x}(\sqrt{x}+1)} \cdot \frac{\sqrt{x}+1}{2} \\
			&= \frac{\sqrt{x} + 2}{\sqrt{x}(\sqrt{x}+1)} \cdot \frac{\sqrt{x}+1}{2}
		\end{align*}
		Vì $x>0$ nên $\sqrt{x}+1 \ne 0$. Ta có thể rút gọn $\sqrt{x}+1$.
		\begin{align*}
			P &= \frac{\sqrt{x} + 2}{\sqrt{x}} \cdot \frac{1}{2} \\
			&= \frac{\sqrt{x} + 2}{2\sqrt{x}}
		\end{align*}
		Vậy $P = \frac{\sqrt{x} + 2}{2\sqrt{x}}$.
	}
\end{ex}

\begin{ex} % Câu 10
	Cho phương trình $x^2-2x-3=0$, không giải phương trình, hãy tính giá trị biểu thức: $A = x_1^2 + x_2^2 - x_1x_2$.
	\loigiai{
		Phương trình đã cho là $x^2-2x-3=0$.
		Đây là phương trình bậc hai có dạng $ax^2+bx+c=0$ với $a=1, b=-2, c=-3$.
		Ta tính biệt thức $\Delta = b^2-4ac = (-2)^2 - 4(1)(-3) = 4 + 12 = 16$.
		Vì $\Delta = 16 > 0$ nên phương trình có hai nghiệm phân biệt $x_1, x_2$.
		
		Theo hệ thức Vi-ét, ta có:
		$$ \begin{cases} 
			x_1+x_2 = \dfrac{-b}{a} = \dfrac{-(-2)}{1} = 2 \\ 
			x_1x_2 = \dfrac{c}{a} = \dfrac{-3}{1} = -3 
		\end{cases} $$
		Ta cần tính giá trị biểu thức $A = x_1^2 + x_2^2 - x_1x_2$.
		Biến đổi $x_1^2 + x_2^2 = (x_1+x_2)^2 - 2x_1x_2$.
		Vậy,
		\begin{align*}
			A &= (x_1+x_2)^2 - 2x_1x_2 - x_1x_2 \\
			&= (x_1+x_2)^2 - 3x_1x_2
		\end{align*}
		Thay giá trị $x_1+x_2 = 2$ và $x_1x_2 = -3$ vào biểu thức $A$:
		\begin{align*}
			A &= (2)^2 - 3(-3) \\
			&= 4 - (-9) \\
			&= 4 + 9 \\
			&= 13
		\end{align*}
		Vậy giá trị của biểu thức $A$ là 13.
	}
\end{ex}

\begin{ex} % Câu 11
	Giải hệ phương trình: $\begin{cases} 2x-3y = 1 \\ 5x - y = 9 \end{cases}$.
	\loigiai{
		Ta có hệ phương trình:
		$$ \begin{cases} 2x-3y = 1 & (1) \\ 5x - y = 9 & (2) \end{cases} $$
		Từ phương trình (2), suy ra $y = 5x - 9$.
		Thay $y = 5x - 9$ vào phương trình (1):
		\begin{align*}
			2x - 3(5x - 9) &= 1 \\
			2x - 15x + 27 &= 1 \\
			-13x &= 1 - 27 \\
			-13x &= -26 \\
			x &= \frac{-26}{-13} \\
			x &= 2
		\end{align*}
		Thay $x=2$ vào $y = 5x - 9$:
		\begin{align*}
			y &= 5(2) - 9 \\
			y &= 10 - 9 \\
			y &= 1
		\end{align*}
		Vậy nghiệm của hệ phương trình là $(x; y) = (2; 1)$.
	}
\end{ex}
\begin{ex} % Câu 12 (Từ ảnh mới)
	Một túi đựng các quả cầu có kích cỡ giống nhau và khối lượng giống hệt nhau, khác nhau về màu trong đó 3 quả màu đỏ, 4 quả màu tím, 5 quả màu vàng, 6 quả màu trắng và 2 quả màu đen. Tính xác suất để lấy được quả không phải màu đỏ.
	\loigiai{
		Tổng số quả cầu trong túi là:
		$N = 3 (\text{đỏ}) + 4 (\text{tím}) + 5 (\text{vàng}) + 6 (\text{trắng}) + 2 (\text{đen}) = 20$ quả.
		
		Ký hiệu các quả cầu:
		\begin{itemize}
			\item Đỏ: Đ$_1$, Đ$_2$, Đ$_3$
			\item Tím: T$_1$, T$_2$, T$_3$, T$_4$
			\item Vàng: V$_1$, V$_2$, V$_3$, V$_4$, V$_5$
			\item Trắng: Tr$_1$, Tr$_2$, Tr$_3$, Tr$_4$, Tr$_5$, Tr$_6$
			\item Đen: Đe$_1$, Đe$_2$
		\end{itemize}
		
		Tập hợp tất cả các kết quả có thể xảy ra là:
		 $\Omega = \{ \text{Đ}_1, \text{Đ}_2, \text{Đ}_3, \text{T}_1, ..., \text{T}_4, \text{V}_1, ..., \text{V}_5, \text{Tr}_1, ..., \text{Tr}_6, \text{Đe}_1, \text{Đe}_2 \}$).
		
		Tập hợp các kết quả thuận lợi cho biến cố A là:\\
		$A = \{ \text{T}_1, \text{T}_2, \text{T}_3, \text{T}_4, \text{V}_1, \text{V}_2, \text{V}_3, \text{V}_4, \text{V}_5, \text{Tr}_1, \text{Tr}_2, \text{Tr}_3, \text{Tr}_4, \text{Tr}_5, \text{Tr}_6, \text{Đe}_1, \text{Đe}_2 \}$.
		Số quả cầu không phải màu đỏ là:
		$n(A) = 4 (\text{tím}) + 5 (\text{vàng}) + 6 (\text{trắng}) + 2 (\text{đen}) = 17$ quả.
		
		Xác suất của biến cố A là:
		$$ P(A) = \frac{\text{Số kết quả thuận lợi cho A}}{\text{Tổng số kết quả có thể}} = \frac{n(A)}{n(\Omega)} = \frac{17}{20} $$
		Vậy xác suất để lấy được quả không phải màu đỏ là $\frac{17}{20}$.
	}
\end{ex}
\begin{ex} % Câu 13
	Một người đầu tháng mua một can xăng đầy hết 270 nghìn đồng, đến cuối tháng đó cũng mua 270 nghìn đồng nhưng còn thiếu 1 lít nữa mới đầy can. Hỏi giá xăng cuối tháng đó là bao nhiêu một lít. Biết giá xăng cuối tháng cao hơn giá xăng đầu tháng là 3 nghìn đồng một lít. (Can có dung tích như nhau).
	\loigiai{
		Gọi giá xăng đầu tháng là $x$ (nghìn đồng/lít).
		Điều kiện: $x > 0$.
		Giá xăng cuối tháng là $x+3$ (nghìn đồng/lít).
		
		Dung tích của can xăng (số lít xăng mua được đầu tháng) là: $\frac{270}{x}$ (lít).
		Số lít xăng mua được cuối tháng là: $\frac{270}{x+3}$ (lít).
		
		Theo đề bài, số lít xăng mua được cuối tháng ít hơn dung tích can 1 lít, nên ta có phương trình:
		$$ \frac{270}{x} - \frac{270}{x+3} = 1 $$
		Quy đồng mẫu số và giải phương trình:
		\begin{align*}
			\frac{270(x+3) - 270x}{x(x+3)} &= \frac{x(x+3)}{x(x+3)} \\
			270(x+3) - 270x &= x(x+3) \\
			270x + 810 - 270x &= x^2 + 3x \\
			810 &= x^2 + 3x \\
			x^2 + 3x - 810 &= 0
		\end{align*}
		Tính biệt thức $\Delta = b^2 - 4ac = 3^2 - 4(1)(-810) = 9 + 3240 = 3249$.
		Vì $\Delta > 0$, phương trình có hai nghiệm phân biệt:
		$\sqrt{\Delta} = \sqrt{3249} = 57$.
		$$ x_1 = \frac{-3 + 57}{2(1)} = \frac{54}{2} = 27 $$
		$$ x_2 = \frac{-3 - 57}{2(1)} = \frac{-60}{2} = -30 $$
		Đối chiếu với điều kiện $x > 0$, ta nhận nghiệm $x=27$.
		Vậy, giá xăng đầu tháng là 27 nghìn đồng/lít.
		Giá xăng cuối tháng là $x+3 = 27+3 = 30$ nghìn đồng/lít.
		
		Kiểm tra lại:
		Dung tích can: $\frac{270}{27} = 10$ lít.
		Số lít mua cuối tháng: $\frac{270}{30} = 9$ lít.
		Chênh lệch: $10 - 9 = 1$ lít (đúng theo đề bài).
		
		Vậy giá xăng cuối tháng đó là 30 nghìn đồng một lít.
	}
\end{ex}
\begin{ex} % Câu 13
	Cho nửa đường tròn $(O;R)$ đường kính $AB$. Tiếp tuyến tại điểm $M$ thuộc nửa đường tròn ($M$ khác $A, B$) cắt các tiếp tuyến tại $A$ và $B$ lần lượt tại các điểm $C, D$.
	\begin{enumerate}
		\item[a)] Chứng minh rằng: Tứ giác $ACMO$ nội tiếp.
		\item[b)] Gọi $E$ là giao điểm của $OC$ và $AM$. Vẽ $MH$ vuông góc với $AB$ tại $H$. Gọi $I$ là giao điểm của $MH$ và $BC$. Chứng minh $E$ là trung điểm của $AM$ và $IE$ song song với $AB$.
	\end{enumerate}
	\loigiai{
		\begin{center}
		\begin{tikzpicture}[line join = round, line cap=round,>=stealth,font=\footnotesize,scale=1]
			% Def
			\def\r{3} 
			\def\a{2}
			
			% Coordinates
			\coordinate (O) at (0,0);
			\coordinate (A) at (-\r,0);
			\coordinate (B) at (\r,0);
			\coordinate (C) at (-\r,\a);
			\coordinate (D) at (\r, \r*\r/\a); % D(r, r^2/a)
			
			\coordinate (M) at ($(C)!(O)!(D)$); % M la hinh chieu cua O len CD
			
			\coordinate (K) at ($(C)!-1!(A)$); % C la trung diem cua AK
			
			\coordinate (H) at ($(A)!(M)!(B)$); % H la hinh chieu cua M len AB 
			
			\coordinate (E) at ($(A)!0.5!(M)$); % E la trung diem cua AM (NEW)
			\coordinate (I) at ($(M)!0.5!(H)$); % I la trung diem cua MH (NEW)
			
			% Draw
			\draw (O) circle (\r); % Ve duong tron (O;r)
			
			\draw (A) -- (B) -- (D) -- (C) -- cycle; % Ve A-B-D-C-A
			\draw (C) -- (K) -- (M); % Ve C-K-M
			\draw (C) -- (O);        % Ve C-O
			\draw (A) -- (M);        % Ve A-M
			\draw (E) -- (I);        % Ve E-I (NEW)
			\draw (O) -- (M) -- (H);        % Ve O-M
			\draw (K) -- (B);        % Ve K-B (NEW)
			
			% Ky hieu goc vuong
			% Can \usetikzlibrary{angles, quotes} o preamble
			\pic[draw, angle radius=3mm]{right angle = C--A--O}; % Goc CAO vuong tai A
			\pic[draw, angle radius=3mm]{right angle = M--H--B}; % Goc MHB vuong tai H
			\pic[draw, angle radius=3mm]{right angle = O--B--D}; % Goc OBD vuong tai B
			
			% Fill 
			\fill (O) circle (1.5pt); 
			\fill (A) circle (1.5pt); 
			\fill (B) circle (1.5pt);
			\fill (C) circle (1.5pt); 
			\fill (D) circle (1.5pt); 
			\fill (M) circle (1.5pt);
			\fill (K) circle (1.5pt); 
			\fill (H) circle (1.5pt);
			\fill (E) circle (1.5pt); % Fill diem E (NEW)
			\fill (I) circle (1.5pt); % Fill diem I (NEW)
			
			% Nodes	
			\node[below left] at (O) {$O$}; 
			\node[below left] at (A) {$A$}; 
			\node[below right] at (B) {$B$};
			\node[above left] at (C) {$C$}; 
			\node[above right] at (D) {$D$}; 
			\node[above] at (M) {$M$}; 
			\node[above right] at (K) {$K$}; 
			\node[below] at (H) {$H$}; 
			\node[above] at (E) {$E$}; % Node cho diem E (NEW), vị trí có thể cần điều chỉnh
			\node[right] at (I) {$I$}; % Node cho diem I (NEW), vị trí có thể cần điều chỉnh
			
			% Ve va danh dau doan thẳng bằng nhau
		\end{tikzpicture}
		
		\end{center}
		}
\end{ex}

\begin{ex} % Câu 14
	\begin{enumerate}
		\item[a)] Một người nông dân có 18 000 000 đồng để làm một hàng rào hình chữ E dọc theo một con sông bao quanh hai khu đất trồng rau có dạng hai hình chữ nhật bằng nhau (hình vẽ). Đối với mặt hàng rào song song với bờ sông thì chi phí làm hàng rào là 120 000 đồng/mét, còn đối với ba mặt hàng rào song song nhau thì chi phí làm hàng rào là 100 000 đồng/mét, mặt giáp với bờ sông không phải rào. Tìm diện tích lớn nhất của hai khu đất thu được sau khi làm hàng rào.
		\item[b)] Cho $a, b$ là các số thực dương thoả mãn $a+b+3ab = 1$.
		Tìm giá trị lớn nhất của biểu thức $P = \sqrt{1-a^2} + \sqrt{1-b^2} + \frac{3ab}{a+b}$.
	\end{enumerate}
\end{ex}



\setcounterref{ex}{0}